%-------------------------------------------------
% FileName: chap-6.tex
% Description: 第6章 系统测试
%-------------------------------------------------

\chapter{系统测试}

\section{测试环境与方法}

为了检验体育馆管理系统是否满足设计要求，本文在前后端联调之后对系统进行综合测试。测试过程主要是对用户端预约业务、器材租借业务、管理端后台维护功能进行测试，测试内容为功能正确性、业务流程完整性、异常处理的有效性、系统运行的稳定性等各方面。根据本课题的实现情况，本文主要用黑盒测试的方法对系统的输入和输出进行对比，来检验系统是否达到了预期的设计目的\cite{zhao2026}。

本次测试所用的环境有前端运行环境、后端服务环境、数据库环境以及浏览器环境这四种。测试时用本地部署的方式运行前后端项目，前端主要实现页面访问和交互，后端主要是处理接口请求以及进行业务校验，数据库用来保存用户、场地、订单、器材等重要的业务数据。具体的测试环境如表 \ref{tab:test-env} 所示。

\begin{table}[htbp]
\centering
\caption{系统测试环境}
\label{tab:test-env}
\small
\begin{tabularx}{\textwidth}{p{3.2cm}p{4cm}X}
\toprule
测试项目 & 环境配置 & 说明 \\
\midrule
操作系统 & Windows 11 & 系统开发与测试使用环境 \\
前端运行环境 & Node.js 18，Vue3，Vite & 用于启动用户端与管理端页面 \\
后端运行环境 & Java 17，Spring Boot 2.7 & 用于提供业务接口与权限控制 \\
数据库环境 & MySQL 8.0 & 用于存储系统业务数据 \\
开发工具 & IntelliJ IDEA，VS Code & 用于前后端代码开发与调试 \\
测试浏览器 & Google Chrome & 用于页面访问与交互测试 \\
\bottomrule
\end{tabularx}
\end{table}

\section{功能测试分析}

根据系统业务流程和功能模块划分，本文选取了用户登录、场地查询、预约下单、订单状态操作、器材租借以及后台管理等关键功能作为测试对象。测试时分别输入正常数据和异常数据，观察系统是否能够返回正确结果，并检查各业务环节之间的数据是否保持一致。部分核心功能测试结果如表 \ref{tab:function-test} 所示。

\begin{table}[t]
\centering
\caption{核心功能测试结果}
\label{tab:function-test}
\small
\begin{tabularx}{\textwidth}{p{1.8cm}p{4.6cm}X p{1.8cm}}
\toprule
测试编号 & 测试内容 & 预期结果 & 测试结果 \\
\midrule
T1 & 输入正确用户名和密码进行登录 & 登录成功并进入系统首页 & 通过 \\
T2 & 按场地类型筛选场地列表 & 返回符合条件的场地数据 & 通过 \\
T3 & 选择未被占用时段提交预约 & 成功生成待支付订单 & 通过 \\
T4 & 选择已被占用时段提交预约 & 系统提示预约冲突并拒绝创建订单 & 通过 \\
T5 & 对待支付订单执行取消操作 & 订单状态更新为已取消 & 通过 \\
T6 & 提交器材租借请求且库存充足 & 生成租借记录并减少可借数量 & 通过 \\
T7 & 提交器材租借请求且库存不足 & 系统提示库存不足并拒绝租借 & 通过 \\
T8 & 管理员修改场地状态和器材信息 & 数据更新成功并在页面中正确显示 & 通过 \\
\bottomrule
\end{tabularx}
\end{table}

从测试结果来看，系统各项核心功能均能够按照预期完成。普通用户可以顺利进行场地浏览、预约下单、订单支付与取消以及器材租借等操作，管理员也能够完成场地维护、订单查看、用户管理和器材信息维护等后台操作。在异常场景下，例如未登录访问受限资源、预约时间冲突、器材库存不足和权限不匹配等情况，系统均能够返回明确提示信息并阻止非法操作，说明系统在业务逻辑控制和异常处理方面具备较好的正确性。

综合分析可知，本系统已经基本实现体育馆管理的主要功能目标，能够满足校园场馆日常预约与管理需求。虽然在大规模并发访问、更加复杂的数据统计和更精细的安全控制方面仍有进一步优化空间，但从当前测试结果来看，系统在功能完整性、流程连贯性和运行稳定性方面均达到了本科毕业设计的预期要求。
